\documentclass[11pt,a4paper]{article}
\usepackage[utf8]{inputenc}
\usepackage[vietnamese]{babel}
\usepackage{amsmath,amssymb}
\usepackage{geometry}
\geometry{a4paper, total={180mm,260mm}, left=15mm, right=15mm, top=18mm, bottom=18mm}

\usepackage{tikz}
\usetikzlibrary{shapes.geometric, arrows.meta, positioning, calc}
\tikzset{
  arr/.style={-{Stealth[scale=0.8]}},
  every node/.append style={align=center}
}
\usepackage{tcolorbox}
\tcbuselibrary{skins,breakable}
\usepackage{fontawesome}
\usepackage{fancyhdr}
\usepackage{enumitem}

\pagestyle{fancy}
\fancyhf{}
\renewcommand{\headrulewidth}{0.5pt}
\renewcommand{\footrulewidth}{0.5pt}
\fancyhead[L]{\small\bfseries\color{violet!85!black}\faBook\ SỔ TAY KIẾN THỨC TOÁN 2}
\fancyhead[R]{\small\bfseries\color{violet!85!black}Lớp Toán Thầy Lĩnh -- loptoanthaylinh.vn}
\fancyfoot[C]{\small\color{gray}-- Sổ tay kiến thức: Bài 35. Ôn tập đo lường --}

\setlength{\parindent}{0pt}
\setlength{\parskip}{0pt}

\begin{document}

% --- BANNER BÀI HỌC ---
\begin{tcolorbox}[
  enhanced,
  colback=violet!8!white,
  colframe=violet!70!black,
  arc=3mm, boxrule=1.2pt,
  left=3mm, right=3mm, top=2mm, bottom=2mm,
  halign=center
]
  {\LARGE\bfseries\color{violet!90!black}\faBookmark\ BÀI 35: ÔN TẬP ĐO LƯỜNG}\\[1mm]
  {\itshape\footnotesize\color{black!75} Ôn thi HK1: Hệ thống 4 đại lượng đo lường: Độ dài, Khối lượng, Dung tích và Thời gian trên 1 trang A4}
\end{tcolorbox}

\vspace{-2mm}

\noindent
\begin{minipage}[t]{0.49\textwidth}\vspace{0pt}
\begin{tcolorbox}[
  enhanced,
  colback=violet!4!white,
  colframe=violet!75!black,
  arc=2mm, boxrule=0.8pt,
  left=2mm, right=2mm, top=1.8mm, bottom=1.8mm,
  title={\footnotesize\bfseries\color{white}\faKey\ 1. Bảng 4 đại lượng đo lường HK1}
]
{\footnotesize
\begin{itemize}\setlength{\itemsep}{1pt}\setlength{\leftskip}{-2.5mm}
  \item \textbf{Độ dài:} Xăng-ti-mét (\textbf{cm}), Đề-xi-mét (\textbf{dm}).\\
  $$1\text{ dm} = 10\text{ cm} \quad (10\text{ cm} = 1\text{ dm})$$
  \item \textbf{Khối lượng:} Ki-lô-gam (\textbf{kg}).\\
  Dùng cân đĩa, cân đồng hồ để đo.
  \item \textbf{Dung tích:} Lít (\textbf{$l$}).\\
  Dùng ca đong, chai lọ để đo sức chứa chất lỏng.
  \item \textbf{Thời gian:} Ngày, giờ, phút, tháng.\\
  $1\text{ ngày} = 24\text{ giờ}$; $1\text{ giờ} = 60\text{ phút}$; $1\text{ tuần} = 7\text{ ngày}$.
\end{itemize}
}
\end{tcolorbox}
\end{minipage}
\hfill
\begin{minipage}[t]{0.49\textwidth}\vspace{0pt}
\begin{tcolorbox}[
  enhanced,
  colback=teal!4!white,
  colframe=teal!75!black,
  arc=2mm, boxrule=0.8pt,
  left=2mm, right=2mm, top=1.8mm, bottom=1.8mm,
  title={\footnotesize\bfseries\color{white}\faSliders\ 2. Kĩ năng tính toán và so sánh}
]
{\footnotesize
\begin{itemize}\setlength{\itemsep}{1pt}\setlength{\leftskip}{-2.5mm}
  \item \textbf{Quy tắc cộng trừ:}\\
  Cùng đơn vị mới được cộng trừ trực tiếp. Kết quả phải kèm tên đơn vị đo.\\
  $15\text{ kg} + 8\text{ kg} = 23\text{ kg}$.\\
  $30\ l - 12\ l = 18\ l$.
  \item \textbf{Đổi đơn vị trước khi so sánh:}\\
  So sánh $2\text{ dm}$ và $18\text{ cm}$:\\
  Đổi $2\text{ dm} = 20\text{ cm} > 18\text{ cm} \implies 2\text{ dm} > 18\text{ cm}$.
  \item \textbf{Ước lượng thực tế phù hợp.}
\end{itemize}
}
\end{tcolorbox}
\end{minipage}

\vspace{-1.5mm}
\begin{tcolorbox}[
  enhanced,
  colback=orange!4!white,
  colframe=orange!80!black,
  arc=2mm, boxrule=0.8pt,
  left=2mm, right=2mm, top=1.5mm, bottom=1.5mm,
  title={\footnotesize\bfseries\color{white}\faSitemap\ 3. Sơ đồ tư duy 4 trụ cột Đo lường Toán 2 HK1}
]
\centering
\begin{tikzpicture}[
  node distance=0.8cm,
  box/.style={rectangle, rounded corners=2pt, draw=violet!80!black, fill=violet!15, text width=2.2cm, align=center, font=\scriptsize\bfseries, inner sep=3pt},
  branch/.style={rectangle, rounded corners=2pt, draw=#1!80!black, fill=#1!10, text width=2.4cm, align=center, font=\scriptsize, inner sep=2.5pt},
  arr/.style={-{Stealth[scale=0.8]}, thick, #1!80!black}
]
  \node[box] (root) {ĐO LƯỜNG\\TOÁN 2 HK1};
  \node[branch=teal, above left=0.3cm and 0.4cm of root] (b1) {\textbf{ĐỘ DÀI}\\cm, dm\\$1\text{ dm} = 10\text{ cm}$};
  \node[branch=orange, above right=0.3cm and 0.4cm of root] (b2) {\textbf{KHỐI LƯỢNG}\\Ki-lô-gam (kg)\\Cân nặng};
  \node[branch=green, below left=0.3cm and 0.4cm of root] (b3) {\textbf{DUNG TÍCH}\\Lít ($l$)\\Sức chứa chất lỏng};
  \node[branch=purple, below right=0.3cm and 0.4cm of root] (b4) {\textbf{THỜI GIAN}\\Ngày, giờ, phút, tháng\\$24\text{h}, 60\text{p}, 7\text{ ngày}$};

  \draw[arr=teal] (root) -- (b1);
  \draw[arr=orange] (root) -- (b2);
  \draw[arr=green] (root) -- (b3);
  \draw[arr=purple] (root) -- (b4);
\end{tikzpicture}
\end{tcolorbox}

\vspace{-1.5mm}
\noindent
\begin{minipage}[t]{0.49\textwidth}\vspace{0pt}
\begin{tcolorbox}[
  enhanced,
  colback=green!4!white,
  colframe=green!60!black,
  arc=2mm, boxrule=0.8pt,
  left=2mm, right=2mm, top=1.8mm, bottom=1.8mm,
  title={\footnotesize\bfseries\color{white}\faCheckCircle\ 4. Ví dụ thực hành chuẩn mực}
]
{\footnotesize
\textbf{1. Điền dấu ($>, <, =$):}\\
- $3\text{ dm} \dots 28\text{ cm}$ ($30\text{ cm} > 28\text{ cm} \implies >$).\\
- $15\ l + 5\ l \dots 20\ l$ ($20\ l = 20\ l \implies =$).\\
- $1\text{ ngày} \dots 20\text{ giờ}$ ($24\text{ giờ} > 20\text{ giờ} \implies >$).
\textbf{2. Bài toán thực tế:}\\
Mẹ mua $5\text{ kg}$ gạo nếp và $18\text{ kg}$ gạo tẻ. Hỏi mẹ mua tất cả bao nhiêu ki-lô-gam gạo?\\
\textit{Giải:} $5 + 18 = \mathbf{23\text{ kg}}$.
}
\end{tcolorbox}
\end{minipage}
\hfill
\begin{minipage}[t]{0.49\textwidth}\vspace{0pt}
\begin{tcolorbox}[
  enhanced,
  colback=red!4!white,
  colframe=red!75!black,
  arc=2mm, boxrule=0.8pt,
  left=2mm, right=2mm, top=1.8mm, bottom=1.8mm,
  title={\footnotesize\bfseries\color{white}\faExclamationTriangle\ 5. Mẹo tránh bẫy trắc nghiệm}
]
{\footnotesize
\begin{itemize}\setlength{\itemsep}{1pt}\setlength{\leftskip}{-2.5mm}
  \item \textbf{Bẫy không đổi đơn vị khi so sánh:}\\
  Thấy $3\text{ dm} < 28\text{ cm}$ vì $3 < 28$ là \textbf{sai}! Phải đổi $3\text{ dm} = 30\text{ cm} > 28\text{ cm}$.
  \item Luôn ghi đơn vị đo ở đáp số của bài toán.
\end{itemize}
}
\end{tcolorbox}
\end{minipage}


\end{document}
